\section{Self-Supervised Pre-Training with DINO}
\label{sec:dino}

% TODO: Discuss the motivation for using self-supervised learning on remote sensing/aerial imagery.
% Highlight the challenges in annotating large-scale poultry farm datasets and why DINO is a suitable foundation model approach.

\subsection{DINO Pre-Training Strategy}
% TODO: Detail the architecture (ViT-Base).
% - Dataset used for pre-training (e.g., NAIP aerial imagery, number of images).
% - Training hyperparameters: 400 epochs, bf16 precision, batch size 32.
% - Multi-crop strategy: 2 global crops (224x224), 8 local crops (96x96).
% - Optimizer settings: Learning rate (0.0005) with warmup, weight decay schedule, teacher momentum.

\subsection{Downstream Fine-Tuning for Segmentation}
% TODO: Explain how the pre-trained DINO backbone was adapted for the downstream segmentation/detection task.
% - Architecture modifications (e.g., attaching segmentation heads).
% - Fine-tuning process and hyperparameters.
% - Comparison of the DINO-pretrained model vs. ImageNet-pretrained baselines.
% - Summarize key results (e.g., improved mIoU or AP on poultry farm detection).

\subsection{Few-Shot Evaluation}
% TODO: Discuss the evaluation of the model in a few-shot setting to measure sample efficiency.
% - Methodology for few-shot sampling (e.g., 5-shot, 10-shot).
% - Comparison with ImageNet baseline models and learning-from-scratch models.
% - Quantitative results and qualitative observations on representation robustness.
